\cleardoublepage

\newrefsection

\chapter{外文翻译}

\ctexset{subsection/number=\arabic{chapter}.\arabic{section}.\arabic{subsection}}

\refstepcounter{section}
\section*{译文题目：VLN-R1 :通过强化学习的视觉语言导航}
\addcontentsline{toc}{section}{\protect\numberline{\thesection}译文题目：VLN-R1 :通过强化学习的视觉语言导航}

{\zihao{3}
\noindent \textbf{原文题目：VLN-R1: Vision-Language Navigation via Reinforcement Fine-Tuning} \\
\textbf{原文来源：} Qi Z, Zhang Z, Yu Y, et al. VLN-R1: Vision-Language Navigation via Reinforcement Fine-Tuning[J]. arXiv preprint arXiv:2506.17221, 2025. \\
}


\vskip 5mm

{
\setlength{\baselineskip}{16pt}\selectfont{
\noindent {\heiti 摘\quad 要\quad }
视觉语言导航（VLN）是具身人工智能（embodied AI）领域的一项核心挑战，它要求智能体（agent）利用自然语言指令在真实世界环境中进行导航。当前基于语言模型的导航系统在离散的拓扑图上运行，这将其路径规划限制在预定义的节点连接上。我们提出了 VLN-R1，这是一个端到端的框架，它利用大型视觉语言模型（LVLM）将第一人称视角的视频流直接转换为连续的导航动作，并采用了受 DeepSeek-R1 启发的、基于 GRPO 的训练方法。为了实现有效的训练，我们首先使用 3D 模拟器（即 Habitat）构建了 VLN-Ego 数据集，并提出了长短期记忆采样（Long-Short Memory Sampling）方法，以平衡历史观测和当前观测。尽管大型语言模型能够监督完整的文本指令，但它们缺乏细粒度的动作级别控制。我们的框架采用了一种两阶段训练方法：a) 监督式微调（SFT），旨在将模型的动作序列文本预测与专家演示对齐；随后是 b) 强化学习微调（RFT），该微调通过一种时间衰减奖励（Time-Decayed Reward, TDR）机制得到增强，该机制能够策略性地对多步未来动作进行加权。实验结果表明，VLN-R1 在 VLN-CE 基准测试中取得了优异的性能。VLN-R1 证明了，通过数据高效、奖励驱动的后训练（post-training），大型视觉语言模型可以驱动具身导航，并增强任务特定的推理能力。
\par}}
\vspace {5mm}

\subsection{ 引言}
具身人工智能（Embodied AI）通过为智能体（agent）赋予第一人称视角的感知（egocentric perception）和物理行动的能力，正在为机器人领域带来革命性的变革。视觉语言导航（Vision-and-language navigation, VLN）是具身AI领域的一项基础性挑战，其要求智能体必须解读自然语言指令，在不熟悉的3D环境中进行导航。这项任务不仅要求语言理解能力，还要求在场景中进行实时决策。

尽管VLN领域取得了显著进展，但当前的研究仍面临挑战。现有方法通常采用具有固定连接性的、基于图的表示方法，这限制了智能体泛化至未见过的或连续环境的能力。虽然有几项研究已将VLN推广到连续空间（VLN-CE），通过以细粒度的电机控制取代离散路点，从而在模拟器中实现自由移动，但这些方法通常需要诸如深度图和导航地图等额外信息，并依赖于像CLIP这样的专门化模型进行视觉-语言对齐，这限制了它们在人-机交互中的泛化能力。

大型语言模型（LLM）在聊天机器人和自动驾驶等领域展现了令人印象深刻的泛化能力，对具身智能领域产生了影响。最近，大型语言模型/大型视觉语言模型（LLMs/LVLMs）已被用于视觉语言导航（VLN）任务，如（左）图1所示。以往的方法使用语言模型作为规划器，将视觉输入转换为文本用于路径规划。然而，这些方法仍然受到这些预定义导航图的限制，这阻碍了在真实世界环境中创建能够整合视觉、语言和行动的、真正意义上的具身智能体。尽管Navid和Uni-Navi构成了重要的进展，但它们对模块化视觉处理流水线的持续依赖，从根本上限制了其可扩展性和泛化能力。

为了解决LVLMs在VLN中的局限性，我们提出了VLN-R1——一个新颖的框架，该框架通过先进的LVLMs（例如，Qwen2-VL）直接处理第一人称视角的视频流，以在连续环境中进行视觉语言导航，如（图1中第3部分）所示。我们首先引入VLN-Ego，这是一个用于在连续导航任务上训练LVLMs的新型数据集。该数据集由Habitat模拟器生成，包含了与相应未来动作预测配对的第一人称视角视频流。针对视频输入处理，我们实现了一种长短期记忆采样（Long-Short Memory Sampling）策略，该策略动态地平衡了历史帧的重要性与实时输入的敏感性。

如图1中第3部分所示，VLN-R1生成四种基本动作指令：前进（FORWARD）、左转（TURN-LEFT）、右转（TURN-RIGHT）、停止（STOP）。训练过程包括两个阶段：a) 监督式微调（Supervised Fine-Tuning, SFT）：在此阶段，模型的文本输出受到完全监督，以与基准真实文本（ground truth text）对齐。b) 强化学习微调（Reinforcement Fine-Tuning, RFT）：在此阶段，我们应用一种受Deepseek-R1启发的强化学习方法，使用群组相对策略优化（Group Relative Policy Optimization, GRPO）来微调模型。我们提出了一种时间衰减奖励（Time-Decayed Reward, TDR）机制，用于评估和验证多步动作预测，从而显著增强了长时程（long-horizon）导航性能。
VLN-R1实现了业界顶尖（state-of-the-art）的导航性能，并支持具身问答（Embodied Question Answering, EQA，结果见补充材料）。我们在Room-to-Room（R2R）和Room-Across-Room（RxR）基准测试的连续版本上评估了其导航能力，证明了我们模型的有效性。本文做出了以下主要贡献：
\begin{enumerate}
    \item[$\bullet$] 我们引入了VLN-Ego数据集，这是一组包含与未来动作预测配对的第一人称视角视频流的集合，专为在连续环境中训练用于视觉语言导航的LVLMs而设计。
    \item[$\bullet$] 我们提出了VLN-R1，一个端到端的具身AI框架，该框架使用大型视觉语言模型（LVLM）处理第一人称视角的视频流，实现了基于动作的实时视觉语言导航。
    \item[$\bullet$] 我们开创性地将GRPO和强化学习微调（RFT）集成到导航任务的LVLMs训练中。此外，我们设计了一种时间衰减奖励（TDR）机制，以专门优化LVLM在导航场景中的性能。 
\end{enumerate}

\subsection{ 相关工作}

视觉语言导航（Vision-and-Language Navigation, VLN）。 这是一项多模态任务，其中智能体（agent）遵循自然语言指令（例如，“左转，然后在沙发旁停下”）在真实或模拟的视觉环境中导航。早期的VLN方法在离散环境中运行，智能体在这些环境中将语言与视觉观测对齐，并在预定义的节点之间进行传送。然而，这些方法在泛化至开放世界场景方面存在困难。随后，VLN-CE基准测试引入了连续环境，支持低级别动作预测（例如，转向角度、步长），并允许在模拟器中自由移动。相关方法包括预测低级别的控制指令，或选择由路点预测器生成的子目标。后续的进展利用了视觉语言预训练模型或任务特定的预训练框架来提升性能。近来，语言模型（LLM）已被集成到VLN系统中。

作为VLN智能体的语言模型。 大型语言模型（LLM）的最新进展已将视觉语言导航（VLN）的研究重点从任务特定的优化转向了开放域的指令泛化。当前的方法利用LLM来提升多任务导航能力和跨场景适应性。它们将现成的（off-the-shelf）LLM与视觉基础模型相结合，以将环境观测转换为文本描述，从而允许语言模型引导智能体。然而，这种方法存在误差累积问题，并且大多局限于离散环境。另一趋势是将大型视觉语言模型（LVLM）集成用于连续导航，但这需要复杂的额外架构。我们的方法在一个类似于Deepseek-R1的框架内，使用基于第一人称视角视频训练的大型视觉语言模型，并采用强化学习微调进行训练。

强化学习微调（Reinforcement Fine-Tuning, RFT）。 大型视觉语言模型（LVLM）的后训练（Post-training）传统上使用监督式微调（SFT），借助带标注的多模态数据来提升任务特定的性能。尽管强化学习微调（RFT）已在语言模型中增强了推理能力（例如，数学问题和代码生成），但其在LVLM上的应用一直有限，主要集中于幻觉缓解或人类偏好对齐。受GRPO和Deepseek-R1-Zero的启发——它们证明了仅凭强化学习（RL）便可提升推理能力——我们开创性地将可验证的奖励机制和基于GRPO的策略优化集成到视觉语言导航任务中。这是RFT在VLN中首次被应用于连续导航决策。

\subsection{ 方法论}
在本节中，我们介绍我们的方法论。第3.1节涵盖了基于强化学习（RL）的后训练（post-training）预备知识以及GRPO。第3.2节描述了我们从Habitat模拟器进行数据收集的流程。第3.3节详细阐述了使用监督式微调（SFT）的第一阶段训练，而第3.4节则描述了第二阶段的强化学习微调（RFT）过程。

\subsubsection{预备知识}
可验证奖励的强化学习（Reinforcement learning with verifiable rewards, RLVR）。 这是一种针对大型语言模型的新型后训练方法，主要面向具有客观基准真相（ground-truth）答案的任务，例如数学和代码生成。相比之下，先前的人类反馈强化学习（Reinforcement Learning from Human Feedback, RLHF）依赖人类指导来学习“什么是正确的”。RLVR不仅具有更简洁的奖励函数，而且能更好地与任务的内在评估标准对齐。具体而言，给定一个输入问题 $ q $，策略模型 $ \pi_\theta $ 会生成响应 $ o $。RLVR的优化目标可以公式化为：

\begin{equation}
\max_{\pi_\theta} \mathbb{E}_{o \sim \pi_\theta(q)} [R_{\text{RLVR}}(q, o)] = R(q, o) - \beta \text{KL} [\pi_\theta(o|q) \parallel \pi_{\text{ref}}(o|q)]
\end{equation}

此处，$ \pi_{\text{ref}} $表示优化前的参考模型。此外，$ \beta $  是控制KL散度（KL-divergence）惩罚比例的超参数。项 $ R(q, o) $ 对应于可验证的奖励函数。具体的目标函数公式化为：

\begin{equation}
R(q, o) =
\begin{cases}
  1, & \text{如果 } o = \text{基准真相}, \\
  0, & \text{其他情况}.
\end{cases}
\end{equation}

在此设定中，$ R(q, o) $ 代表一个硬奖励函数（hard reward function），即当输出 $ o $ 与基准真相匹配时，奖励设为1，否则为0。与此不同，我们的方法引入了软奖励函数（soft reward functions），以实现更精细的优化目标。

R1-Zero与GRPO。 DeepSeek R1-Zero算法使用群组相对策略优化（Group Relative Policy Optimization, GRPO）进行强化学习。与需要一个评论家模型（critic model）的PPO或依赖离线偏好数据的DPO不同，GRPO在线运行且无需一个显式的奖励模型，仅使用群组范围的奖励归一化（group-wise reward normalization），这使其非常适用于大型语言模型。对于一个问题 $ q $，GRPO从策略模型 $ \pi_\theta $ 生成 $ G $ 个响应 $ {o_1, \dots, o_G} $ 并计算出奖励 $ {r_1, \dots, r_G} $。然后，它对这些奖励进行归一化以确定相对优势：

\begin{equation}
A_i = \frac{r_i - \mu\{r_1, \dots, r_G\}}{\sigma\{r_1, \dots, r_G\}},
\end{equation}

其中 $ A_i $ 代表第 $ i $ 个回答的相对质量。GRPO鼓励模型偏好组内奖励值更高的响应。GRPO的优化目标是：

\begin{equation}
J_{\text{GRPO}}(\theta) = \mathbb{E} \left[ \sum_{i=1}^{G} \min\left( \frac{\pi_\theta(o_i|q)}{\pi_{\text{ref}}(o_i|q)} A_i, \text{clip}\left(\frac{\pi_\theta(o_i|q)}{\pi_{\text{ref}}(o_i|q)}, 1 - \epsilon, 1 + \epsilon \right) A_i \right) \right] - \beta D_{\text{KL}}(\pi_\theta \parallel \pi_{\text{ref}})
\end{equation}

比率 $ \frac{\pi_\theta(o_i|q)}{\pi_{\text{ref}}(o_i|q)} $ 量化了偏好转移（“>1”表示对 $ o_i $ 有更高的偏好）。优势 $ A_i $对更新进行缩放，裁剪（clipping）和最小化（min）操作确保修改的保守性：适度提升高奖励响应（$ A_i $> 0），同时谨慎地惩罚低劣响应（$A_i$ < 0）。KL散度可防止模型对 $ \pi_{\text{ref}} $ 产生过度偏离。

\begin{equation}
D_{\text{KL}} = \pi_{\text{ref}}(o_i|q) \log\left( \frac{\pi_{\text{ref}}(o_i|q)}{\pi_\theta(o_i|q)} \right)
\end{equation}

\subsubsection{数据集引擎：VLN-Ego}

对于室内导航任务，先前的非LVLM方法要么在离散环境中运行，要么依赖于诸如全局地图和深度信息等额外模态。与先前的工作相比，我们的框架使基于LVLM的导航能够仅使用第一人称视角（egocentric）的视觉输入，从而消除了对额外传感器模态的依赖。这使得LVLM智能体能够在一个连续的环境中运行，在其中它可以执行灵活的动作，例如按特定距离前进或按特定角度转向。因此，我们设计了 VLN-Ego，一个基于第一人称视角视频流的、用于导航的LVLM训练引擎。

我们利用了Meta的Habitat模拟平台，该平台构建于Matterport3D 场景数据集之上。总共90个场景被划分为训练集（61个场景）、可见验证集（11个场景）和未见验证集（18个场景），分别用作训练、验证和测试。为了与先前的基准保持一致性，我们还在Matterport3D场景中采用了R2R(Room-to-Room) 和 RxR  (Room-across-Room) 的轨迹。R2R主要涉及从一个房间导航到另一个房间，包含7,189条路径，其指令例如：\textit{“走过厨房餐桌，在走廊左转，然后在卧室门口停下。”} 另一方面，RxR的路径长度是R2R的三倍，我们选择了其英语语言指令，总计42,023条轨迹。

VLN-Ego 逐个动作地（action-by-action）获取标注，其标注结构详见图2。它主要由三部分构成。第一部分是指令部分，包括 \texttt{<System Message>} 和 \texttt{<Instruction>}。后者指的是作为输入的自然语言指令。

\subsubsection{监督式微调 (SFT) 阶段}


VLN-R1 任务被正式定义如下：在每个时间步 $ t $，给定一条自然语言指令 $ \mathcal{I} $ 和一个由历史观测组成的第一人称视角视频流 $ V_{0:t-1} = \{v_0, v_1, \dots, v_{t-1}\} $，智能体需要预测一个包含 $ n $ 个未来动作的序列 $ \{A_t, A_{t+1}, \dots, A_{t+n-1}\} $，以便在新的环境中成功导航。每个动作 $ A_t $ 对应一个低级别的原子运动。动作空间由四个基本动作构成：$ \{\text{前进, 左转, 右转, 停止}\} $。如果智能体在目标的阈值距离内停止，则认为该导航片段成功。


传统的历史帧选择策略存在局限性。均匀采样对所有过去帧赋予同等权重，忽略了近期观测的相关性，而指数衰减则过度优先考虑邻近帧，从而丢失了在导航中至关重要的长期上下文。为克服这些局限性，我们提出了一种长短期记忆采样（Long-Short Memory Sampling）策略，旨在平衡短期相关性与长期上下文感知。其形式化定义如下：

\begin{equation}
H_t = \underbrace{\{v_{t-1 \cdot \delta_1}, v_{t-2 \cdot \delta_1}, \dots, v_{t-M}\}}_{\text{短期记忆 (采样率 } \delta_1 \text{)}} \cup \underbrace{\{v_{t-M-\delta_2}, v_{t-M-2\delta_2}, \dots, v_0\}}_{\text{长期记忆 (采样率 } \delta_2 \text{)}}
\end{equation}

短期记忆在 $ M $ 步内每 $ \delta_1 $ 帧进行一次采样，而长期记忆则以 $ \delta_2 $ 的间隔（$ \delta_2 > \delta_1 $）对余下的历史进行采样。这种方法既保留了近期的细节，也保留了长距离的上下文。



\subsubsection{强化学习微调 (RFT) 阶段}

在监督式微调（SFT）之后，我们在第二阶段使用强化学习微调（RFT）对模型进行进一步优化。我们沿用相同的指令数据格式，但设计了一种创新的奖励函数。以往的强化学习微调采用简化的奖励机制，例如二元选项匹配、误差阈值或诸如 ROUGE-L 等文本相似度指标。然而，这些方法对于视觉语言导航（VLN）任务而言是不充分的，因为它们缺乏对时序预测的有效监督。

我们提出一种时间衰减奖励（Time-Decayed Reward, TDR）函数来处理时序依赖性。我们的方法从预测的文本元组 $ (a_{t+k}, \phi(a_{t+k})) $ 中提取动作分量 $ a_{t+k} $，以形成未来 $ n $ 步的序列化预测：$ A_{t:t+n} = (a_t, a_{t+1}, \dots, a_{t+n-1}) $。我们的时序加权方案将指数衰减 $ \gamma $ 应用于动作的重要性，确保了近端动作比远端动作获得显著更高的权重。该奖励函数如下所示：

\begin{equation}
R_{\text{nav}} = \sum_{k=0}^{n-1} \gamma^k \cdot \mathbb{I}(a_{t+k} = a_{t+k}^*)
\end{equation}

其中 $ \mathbb{I} $ 是指示函数，当预测动作与基准真相 $ a_{t+k}^* $ 匹配时，其输出为1，否则为0。该设计通过应用指数衰减项 $ \gamma^k $ 来确保位置感知的（position-aware）一致性，该项优先考虑更早的正确动作。我们的方法克服了监督式微调（SFT）的局限性，SFT只能优化整体的文本输出，而无法为动作序列的生成提供精确的指导。

\subsection{ 实验}

我们在此呈现我们的实验设置与结果。在第4.1节中，我们解释了实现细节。第4.2节展示了我们的主要结果，而第4.3节则进行了消融研究。

\subsubsection{实现细节}

我们利用 VLN-Ego 作为我们的训练数据集，在 SFT 阶段使用了来自 R2R 和 RxR 数据集的共计 1.8M 样本。在 RFT 阶段，我们从每个数据集中随机选取 10K 样本，共计 20K 样本用于训练。智能体仅使用第一人称视角的视频来导航至目标，如果智能体在阈值距离内停止，则导航片段被视为成功。我们在18个未见过的场景（Val-Unseen）上，使用 VLN-CE 的指标进行评估。SR$\uparrow$（成功率）、OS$\uparrow$（谕示成功率）和 SPL$\uparrow$（成功率按路径长度加权）用于衡量导航准确性，而 NE$\downarrow$（导航误差）和 TL$\downarrow$（轨迹长度）用于评估效率。所有距离测量单位均为米（m）。我们训练了 Qwen2-VL-2B 和 Qwen2-VL-7B 两个模型，其中 7B 模型部署在8块 NVIDIA A800 GPU 上，并使用 DeepSpeed ZeRO-3 进行优化。在训练期间，图像被缩放到最大分辨率为 65,536 像素，并以每个实例16张图像的批次进行处理（生成 4.1K 词元）。对于监督式微调（SFT），我们采用了 5e-6 的学习率，并结合余弦调度（10\% 的预热），每个 GPU 的批处理大小为2，通过4次梯度累积步骤实现了64的全局批处理大小，在36小时内完成1个周期的训练。对于强化学习微调（RFT），我们将学习率降低至 1e-6，权重衰减为0.01，并设置 $ \beta=0.04 $，将 GRPO 配置为每个提示生成8个样本，使用每个 GPU 的批处理大小为1（无梯度累积），并在大约12小时内完成训练。

\subsubsection{主要结果}

我们的主要实验包括了在 VLN-CE 的 R2R 和 RxR 上进行训练的结果。我们在 Val-Unseen 划分上进行评估，该划分包含了训练期间完全未见过的环境。除了展示我们方法的有效性，我们还进一步探究了 RFT 在此框架中的作用。

我们在 VLN-CE R2R 上的结果如表1所示。首先，我们的方法仅依赖 RGB 视频输入，同时实现了业界顶尖的性能。通过 RFT 训练的 2B 模型性能优于 SFT 训练的 7B 模型，这与 Deepseek-R1 的发现一致，即强化学习微调（RFT）能使较小模型达到与较大模型相当的性能。

在 VLN-CE RxR 上的结果如表2所示。最后四行展示了使用完整数据集进行训练的结果，其结论与在 R2R 上的结论一致。中间四行代表一种配置，其中监督式微调（SFT）仅使用了 630K R2R 样本，而强化学习微调（RFT）则使用了 10K RxR 样本。令人惊讶的是，仅通过 10K 样本进行 RFT 微调的 LVLM，其性能超越了在完整数据集上训练的对应模型。值得注意的是，RFT 实现了有效的跨领域适应：在一个领域进行基础训练后，仅需一小部分新领域的数据就足以实现卓越的性能迁移。我们在图4中展示了定性的实验结果，并在附录中提供了更多的定性示例和具身问答（Embodied QA）结果。

\subsubsection{消融研究}

我们进行了一系列消融实验，以分析模型在监督式微调（SFT）和强化学习微调（RFT）阶段的性能。所有实验均使用仅在 R2R 上训练的 Qwen2-VL-7B 模型。

\textbf{监督式微调（SFT）阶段。} 我们研究了动作空间（Action Space）和历史记忆（History Memory）组件的不同设计如何影响 VLN-R1 的性能。在这些实验中，我们仅在训练期间使用监督式微调（SFT）。如表3a所示，我们的分析表明，连续预测未来6个动作能产生最优结果。我们观察到，当仅预测单个动作时，性能会出现显著下降，因为这种方法未能考虑未来步骤的依赖关系。表3b证明，我们提出的长短期记忆（Long-Short Memory）策略通过有效平衡当前观测与历史上下文，实现了最佳的帧选择性能。

\textbf{强化学习微调（RFT）阶段。} 我们进一步研究了强化学习微调（RFT）中的超参数。首先，我们在表3c中检验了生成次数（number of generations）的影响，随后探讨了不同奖励函数及其对结果的影响。我们发现，从6代增加到8代的性能提升是微乎其微的，因此我们选择8代。表3d展示了不同奖励函数的效果，证明了我们提出的时间衰减奖励（Time-Decayed Reward）函数的有效性。


\subsection{结论与局限性}

我们提出了 VLN-R1，这是一个端到端的框架，它利用大型视觉语言模型（LVLM）来执行连续的视觉语言导航（VLN）任务。我们的框架通过LVLM处理第一人称视角的视觉输入，无需依赖导航图或额外的传感器。VLN-R1集成了强化学习微调（RFT）、群组相对策略优化（GRPO）和时间衰减奖励（TDR），提升了长时程决策能力。我们的方法在VLN-CE上达到了业界顶尖（SOTA）水平，其中2B模型在经过RFT后性能可与7B模型相媲美，并展示了仅需少量数据即可实现高效的跨领域适应能力。

我们的方法存在局限性，包括评估仅在模拟的室内环境中进行，这限制了其真实世界的泛化能力。离散的动作空间也限制了精细化的控制。尽管如此，我们的工作通过集成Qwen2-VL，将VLN确立为LVLM的一项下游任务，从而推动了导航领域的发展。


\refstepcounter{section}
\section*{译文题目：OmniVLA 全模态机器人导航模型}
\addcontentsline{toc}{section}{\protect\numberline{\thesection}译文题目：OmniVLA 全模态机器人导航模型}

{\zihao{3}
\noindent \textbf{原文题目：OmniVLA: An Omni-Modal Vision-Language-Action Model for Robot Navigation} \\
\textbf{原文来源：} Hirose N, Glossop C, Shah D, et al. OmniVLA: An Omni-Modal Vision-Language-Action Model for Robot Navigation[J]. arXiv preprint arXiv:2509.19480, 2025. \\
}

\vskip 5mm

 {
\setlength{\baselineskip}{16pt}\selectfont{
\noindent {\heiti 摘\quad 要\quad }
人类在导航至目的地时，能够灵活地解读和组合不同的目标指令，例如语言指令、空间坐标或视觉参考。相比之下，大多数现有的机器人导航策略仅针对单一模态进行训练，这限制了它们在现实世界场景中的适应性，而在这些场景中，不同形式的目标指令既自然又互补。在本文中，我们提出了一个用于机器人基础模型的训练框架，该框架能够实现基于视觉导航的全模态（omni-modal）目标条件化。我们的方法利用一个高容量的视觉-语言-动作（VLA）骨干网络，并通过一种随机化的模态融合策略，使用三种主要的目标模态进行训练：2D 位姿、自目心图像和自然语言，以及它们的组合。这种设计不仅扩展了可用数据集的范围，还促使策略发展出更丰富的几何、语义和视觉表征。由此产生的模型 OmniVLA，在未见过的环境中实现了强大的泛化能力，对稀疏模态具有鲁棒性，并能遵循新颖的自然语言指令。我们证明了 OmniVLA 在所有模态上均优于专门的基线模型，并为微调至新模态和新任务提供了一个灵活的基础。我们相信，OmniVLA 为实现广泛通用和灵活的导航策略迈出了一步，并为构建全模态机器人基础模型提供了一条可扩展的路径。我们提供了展示 OmniVLA 性能的视频，并将在我们的项目页面\footnote{\url{omnivla-nav.github.io}}上发布其检查点和训练代码。
\par}}

\subsection{引言}

在环境中导航时，人类会自然地使用多种信息模态来推断和发现通往目标的路径（例如，在地图上查找 GPS 位置、利用视觉地标以及遵循语言指示）。然而，先前在机器人导航领域的研究通常基于狭窄的应用场景，使用单一模态来训练策略。当目标在附近时，用语言来描述它们很方便（例如，“沿着大楼走，然后去入口”），而更远的目标则可以更有效地用 GPS 坐标来描述。然而，一个真正通用的导航策略必须能够执行需要我们利用多种信息来源才能成功的任务。例如，对于自主配送或巡检等任务，指定 GPS 坐标和地标图像的组合可能是最有效的方式。尽管这些模态在很大程度上存在重叠，但它们为任务及其执行方式提供了补充信息，尤其是在机器人从其传感器接收到的世界部分视图的背景下。

受此需求驱动，本研究旨在训练一个能够根据以多种模态表达的目标指令进行导航的通用策略。通过在全模态目标上进行训练，我们旨在实现更强大、更灵活的策略，最终获得一个对新模态和未见环境表现出高度适应性的基础模型。

训练基础模型需要我们尽可能多地利用数据。随着机器人社区近来在数据收集方面的努力，训练强大的通用导航策略已变得可行。然而，这些策略通常仅使用单一类型的任务表征进行训练，例如自目心图像、2D 位姿或自然语言。这限制了可用的数据集，使其仅限于那些符合期望任务表征的数据集（例如，仅带有语言标签的数据集），也限制了模型在测试时的使用方式，并可能限制模型处理任务和观测的方式——即，更偏向空间或视觉的任务可能会提升模型的空间推理能力，而语言任务则可能在增强语义理解方面提供补充优势，这与多模态语言模型的任务混合非常相似。

在本研究中，我们提出了一系列用于自主导航的全模态视觉-语言-动作模型（Omni-Modal Vision-Language-Action Models, OmniVLA），该模型可以接纳以多种模态表达的目标，跨模态利用信息，并实现一个更灵活的导航策略。我们使用通过三种主要模态指定的目标来训练我们的模型：（1）2D 位姿，（2）自目心图像，以及（3）自然语言。通过同时学习解读这些不同的模态，模型必须对任务的几何、视觉和语义信息形成更丰富的理解，从而产生一个更强大的导航模型。

\subsection{相关工作}

在导航任务中，我们可以用多种模态来指定目标，例如自目心图像、2D 位姿和语言。对于以自目心图像为条件的导航，先前的工作整合了各种机器人实体的公开可用数据集来训练通用策略。这些策略在室内环境中表现良好，这些环境需要丰富的视觉信息且无法可靠地接入 GPS。以 2D 位姿为条件的策略在具有 GPS 定位的室外环境中的长时程任务中最为成功。MBRA 引入了一种基于模型的重新标注方法，以利用大规模数据源来执行更具挑战性的、以 2D 目标位姿为条件的远程导航任务。

以语言为条件的导航提供了一种用户友好且灵活的接口，用于指示机器人在给定环境中到达特定的目标物体或区域，即使距离很远。然而，要实现鲁棒的指令遵循，需要的语言理解能力需超越简单的物体指代。早期的工作依赖于预训练的语言编码器，而近期的研究则使用强大的 VLM 骨干网络，或者直接用它们来执行导航任务，或者在机器人数据上对其进行微调。其他工作已扩展到使用反事实动作生成和非机器人数据来改进这些模型的训练。LeLaN 同时利用机器人和非机器人数据来学习一个通用的语言条件导航策略，它采用一种基于模型的方法来生成反事实动作，这些动作能够到达目标物体，并配以源自 VLM 推理的语言提示。尽管这类合成的动作指令和语言提示能够实现可扩展的训练，但它们的不准确性可能成为性能瓶颈。

在机械臂操控领域，近期的机器人基础模型 (RFMs) 旨在统一视觉、语言和动作用于泛化。从早期的多模态 Transformer 发展到专注于灵巧性和真实世界部署的大规模系统，它们标志着向能够实现鲁棒、多功能机器人控制的可扩展框架的转变。一些机械臂操控方面的工作利用了包含不同观测和动作空间以及条件化（语言、位姿等）的大型数据集，通过在训练期间掩码不可用的输入，并展示了在多种输入类型上同时训练的益处。

受先前在机械臂操控领域成功的启发，OmniVLA 超越了以往的导航工作，旨在尽可能多地利用跨多种模态和机器人实体的机器人与非机器人数据，学习一个能够以目标图像、语言和 GPS 为条件的策略。值得注意的是，OmniVLA 在近 10,000 小时的真实世界导航数据上进行训练，据我们所知，这是用于端到端导航策略的最大的预训练数据集。

\subsection{用于导航的全模态 VLA}

我们提出了 OmniVLA，一个具有全模态条件化的端到端导航策略。OmniVLA 能有效捕捉核心导航能力，例如避障和路径跟随，并能通过全模态用户输入执行复杂的目标条件导航行为。此外，为了训练 OmniVLA，我们利用了大量的公开数据集——每个数据集至少包含一种模态——从而得到了一个具有强大泛化能力以及对新模态和新环境适应性更强的模型。

\subsubsection{OmniVLA 架构}

我们在一个高容量的 VLA 架构之上构建 OmniVLA，该架构包含了来自互联网规模预训练的知识（用于基础 VLM）和跨机器人实体的动作数据（用于基础 VLA）。我们修改了该架构以实现灵活的目标条件化和全模态表征学习，同时保留了基础模型中强大的视觉-语言先验。

图 2 展示了该网络架构，它构建于 OpenVLA 之上，一个 7B 参数的 VLA 模型。我们使用一个视觉编码器来处理机器人当前的视觉观测。对于目标条件化，我们的架构支持三种不同的模态——视觉、位置和语言——这些模态也可以同时指定。我们将每个独立的目标模态投射到一个共享的词元空间（token space），作为 LLM 骨干网络的输入。我们使用“模态丢弃”（modality dropout）在训练和推理期间灵活地掩码目标模态。对于动作输出，我们遵循 OpenVLA-OFT 的做法，在 LLM 的输出端添加一个线性动作头，以生成一个动作序列 $\{a_i\}_{i=1...N}$。

我们还实现了一个更小的“边缘版”（edge）OmniVLA 架构，它构建于 ViNT 之上，一个 50M 参数的导航 Transformer。我们对其进行修改，添加了多个目标编码器，并以与我们基础模型相同的方式执行模态丢弃（详见附录 A）。在真实世界评估（第四节）中，我们发现对于资源受限的部署场景，OmniVLA-edge 是一个非常有吸引力的选择，在这些场景中，大型 VLA 模型的推理是不可行的。

\subsubsection{训练 OmniVLA}

尽管使用多模态输入很有吸引力，但训练能够接纳全模态输入的策略需要整合支持训练的机器人数据集，并解决可用模态之间的相对不平衡和稀缺性问题。

\textbf{训练数据。} 表 I 列出了用于视觉导航的公开可用数据集。这些数据由 13 个公开可用的数据集组成，涵盖了 10 种不同机器人实体的 9,500 小时数据。虽然这些数据集中大部分包含人工收集的标签，但我们也纳入了带有合成标签的数据（可能噪声更大）以及通过重新标注过程提炼的动作。据我们所知，我们为训练 OmniVLA 整合了最大规模的导航数据集组合，其中包含了高度多样化的环境、机器人实体和模态。

尽管大型数据集有助于泛化，但大规模的数据收集工作可能会引入更多噪声，从而降低数据准确性。遵循先前的工作，我们在训练中对 Frodobots 数据集使用了由 MBRA 生成的合成动作，并对 LeLaN 数据集使用了 NoMaD 生成的合成动作。由于现有的重新标注方法无法解决 BDD-V 数据集（一个自动驾驶车辆数据集）与我们使用的其他小型机器人数据集之间巨大的实体差异，我们训练了一个重新标注模型来生成合理的合成动作，从而使得利用该数据进行训练成为可能，这与先前的工作类似。然而，我们发现在 GNM 混合数据集中，直接使用高质量的原始动作进行训练是合理的。关于 BDD-V 数据集的详细信息，请参见附录 B。

\textbf{训练方法。} 遵循先前在机械臂操控 VLA 上的工作，我们采用一种随机丢弃的形式在所有可用模态上进行训练，从而得到一个更高效和更通用的模型。为了训练 OmniVLA，我们从表 I 所含的数据集中构建一个批次的样本。对于每个样本，我们再从其可用的目标模态中独立采样，以构成条件输入，我们称之为 $t_m$。例如，对于 GNM 混合数据集，条件输入可以从 2D 位姿、自目心目标图像或它们的组合中选择。自然地，通过使用这种丢弃机制，我们在覆盖所有模态和数据集的同时，也提高了训练的稳定性。

在每个训练步骤中，我们构建一个注意力掩码，该掩码会排除未使用或不可用的模态（用随机值填充），从而确保模型只关注被选中的模态。在这些混合模态的批次上进行训练，可以促使模型更好地表征目标信息，从而为泛化和微调产生更优的表征。

我们训练 OmniVLA 策略，该策略可描述为 $\{\hat{a}_i\}_{i=1...N} = \pi_\theta(I_c, I_g, p_g, l_g, t_m)$，其中 $I_c, I_g, p_g, l_g$ 和 $t_m$ 分别是自目心当前图像、自目心目标图像、2D 目标位姿、语言提示和随机选择的模态。我们通过 $J(\theta) = \sum_{i=1}^{N} (\hat{a}_i - a_i^{\text{ref}})^2$ 来计算目标函数 $J$，以模仿 N 步动作参考 $\{a_i^{\text{ref}}\}_{i=1...N}$，并更新 $\theta$ 以最小化 $J$。遵循之前的工作，LeLaN 数据集中的样本使用一个额外的任务特定目标来鼓励到达物体的行为，其中最终的动作被训练为靠近目标物体。目标函数的详细信息在附录 E 中展示。

\textbf{训练细节。} 我们将动作块大小 N 设为 8，频率为 3 Hz，对应于所有模型 2.4 秒的时长。对于每个批次，我们以 LeLaN : GNM : Frodobots : BDD-V = 4:1:1:1 的比例采样数据以平衡不同模态。由于即使在拥有多个 GPU 的服务器上，我们也无法为某些模型确保足够大的批次大小，因此我们累积多个步骤的梯度以稳定训练过程。在使用 OpenVLA 检查点和八个 H100 GPU 训练 OmniVLA 时，我们使用每个 GPU 7 的批次大小，并累积 4 个步骤的梯度，从而得到一个 224 (=7×8×4) 的有效批次大小。此外，我们应用 LoRA 将可学习参数限制在约 5\%，这使我们能够最大化批次大小，并在训练速度和稳定性之间取得平衡。请注意，由于模型规模较大，LoRA 仅用于基于 OpenVLA 的模型。其他训练设置，如学习率、语言词元化、归一化等，与各类模型原始代码中的默认设置保持一致。

\subsection{实验设置}

我们首先描述我们在真实世界机器人平台上评估全模态导航的实验设置。我们进行了广泛的真实世界评估，并与最先进的专用基线和通用基线进行比较。

\subsubsection{导航任务}

我们考虑以下三种导航任务：

\textbf{以语言为条件的导航：} 我们在语言提示上对 OmniVLA 进行了评估，这些提示不仅将机器人引向目标位置，还具体说明了其沿途应如何行动。我们在 40 个环境中进行了评估，包括办公室、厨房区、入口大厅、公园和人行道，并使用了多样化的语言提示。目标设置在距机器人初始位置 5 到 30 米处。为了评估大型预训练模型的优势，我们引入了分布外（OOD）的语言提示，这些提示超出了训练数据中的指令范围。训练数据主要包含如“朝 X 移动”之类的物体到达指令，而我们设计的 OOD 提示则在一半的试验中明确了导航至目标的方式。遵循 CAST 的方法，我们精心设计了这些提示，并根据机器人对指令的遵守情况来评估其行为。我们还选择了 17 个在机器人起点和目标之间放置了障碍物的环境，使任务更具挑战性，并测试策略的核心导航能力。

\textbf{以自目心目标图像为条件的导航：} 对于自目心目标图像，我们的策略的任务是导航至最远 3 米外的目标位置。为了扩展此范围，我们遵循先前的基于视觉的方法，并采用拓扑记忆在 8 个不同环境中进行导航，从而能够导航至更远的目标。为构建目标图，我们以 1 Hz 的频率记录图像观测。在部署期间，我们从第一个观测开始初始化，并在每个时间步，如先前的工作一样，估计最近的节点作为当前位置。然后，下一个节点的图像作为目标图像 $I_g$ 提供给我们的策略。

\textbf{以目标位姿为条件的导航：} 当以 2D 目标位姿为条件时，我们的策略导航至距离初始位置 25 到 100 米的目标。我们使用 GPS 来估计机器人的位置和目标的位置。在每个步骤中，我们计算机器人相对于目标位姿的位姿，即 $p_g$。在评估中，我们选择了 7 个环境，并在每个环境中于不同时间进行了 3 次试验，以考虑 GPS 的抖动。

\subsubsection{移动机器人平台}

我们在 FrodoBots ERZ 平台上评估 OmniVLA，这是一个低成本的移动机器人，如图 3 所示。它配备了多种传感器，包括前后摄像头、GPS、一个 IMU 单元（陀螺仪、加速度计和指南针）以及所有四个轮子上的轮速传感器。所有的传感器数据流都可以通过一个 Web API 访问，我们训练的导航策略通过发送线性和角速度指令与机器人交互。

为评估跨机器人实体的泛化能力，我们还在另外两个平台上评估了我们的策略：VizBot 轮式机器人和 Unitree Go1 四足机器人。

\subsubsection{基线模型}

在我们的评估中，我们将 OmniVLA 与七个跨所有模态的基线模型进行了比较。这包括从零开始训练的无模型导航策略、能够利用互联网规模视频的方法、使用如 CLIP 等现成视觉表征的方法，以及最先进的 VLA 模型。对于 LeLaN、CounterfactualVLA、ViNT、MBRA-pose 和 MBRA-image，我们使用作者原始的实现和检查点进行评估。下面，我们描述两个与原始实现略有不同的基线模型。

\textbf{NoMaD：} 对于以 2D 目标位姿为条件的导航，我们在探索模式下运行 NoMaD 策略，以生成 30 个候选轨迹。然后我们选择其最终预测位置最接近目标位姿 $p_g$ 的轨迹，并用它来控制机器人。对于以自目心目标图像为条件的导航，我们遵循原始的 NoMaD 实现，使用目标图像 $I_g$。

\textbf{CoW：} 对于以语言为条件的导航，我们将当前观测和描述目标物体的提示提供给 OWL-VIT B/32 检测器——在原始论文中被报告为最强的模型——来估计物体的边界框。遵循之前的方法，我们裁剪边界框内的点云，并计算中值点作为目标物体位姿。为确保与我们的方法（仅依赖单个 RGB 摄像头，无深度或 LiDAR）进行公平比较，我们使用 Depth360 估计深度并重建点云。然后使用一个状态格运动规划器来生成速度指令。

\subsection{评估全模态导航}

为评估我们的 OmniVLA 策略，我们着重通过实验回答以下问题：

\begin{itemize}
    \item[\textbf{Q1}] 全模态预训练是否优于单模态导航策略？
    \item[\textbf{Q2}] OmniVLA 能否遵循多种目标模态的组合？
    \item[\textbf{Q3}] OmniVLA 能否适应新的目标模态、环境和机器人实体？
\end{itemize}

\subsubsection{OmniVLA 与单模态策略的比较}

对于 Q1，我们与单模态策略基线进行了比较分析，结果总结在表 II 中。评估指标在表格标题中提供。我们最大的模型 OmniVLA，在每个特定模态的基线上都显示出明显优势。这种改进源于我们的策略能够从我们庞大且高度多样化的训练数据组合中学习到更通用的导航能力，这个数据组合远大于单个模态可用的数据集。此外，我们观察到，预训练的 VLA 架构和预训练数据的选择对性能有显著影响。值得注意的是，OmniVLA 在以语言和位姿为条件的任务上表现出强大的泛化能力，优于最强的单模态专用模型。

\textbf{架构比较。} 尽管所有 OmniVLA 变体都展示了跨模态的泛化能力，但 OmniVLA 仍然以巨大优势超越了较小的 VLA 模型。我们注意到 MiniVLA 和 SmolVLA 的表现相当差，这主要归因于它们的预训练领域（MiniVLA 主要在机械臂操控数据上训练）与当前任务之间的巨大差距，以及有限的模型容量（SmolVLA 可能没有足够的能力来学习有用的跨机器人实体表征）。相比之下，OmniVLA-edge 采用了更专门针对导航任务的架构（在附录 A 中描述），这使其在同等规模下能够取得卓越的性能。然而，在语言遵循方面，OmniVLA-edge 和 OmniVLA 之间仍存在巨大差距，这凸显了从预训练 VLM 继承的视觉语言先验所带来的好处。

\textbf{图像条件下的性能。} OmniVLA 的性能与最强的单模态基线 MBRA 相当，实现了 100\% 的成功率，并优于 NoMaD 和 ViNT。虽然 NoMaD 和 ViNT 是在更有限的 GNM 混合数据集上训练的，但 MBRA 利用了一个更多样化的预训练数据集，使其能够实现强大的通用性能。OmniVLA 在这些数据源的组合上进行训练，并能够达到同等的强大性能。

\textbf{位姿条件下的性能。} OmniVLA 在成功率（二元）和朝向目标的局部进展方面，相较于最强的专用基线 MBRA-pose，实现了 9\% 的性能提升。尽管两种策略都利用了相同的大型数据集进行位姿条件导航，但凭借更高容量的模型和额外的数据，OmniVLA 能够超越先前的方法。这一观察对于基于 VLA 的策略尤其令人印象深刻，因为这在 VLM/VLA 预训练中并未得到很好的体现。

\textbf{语言条件下的性能。} 在我们评估的多样化语言提示集上，OmniVLA 在目标到达和语言遵循方面都显示出巨大提升。虽然 OmniVLA 和 LeLaN 都仅使用来自 LeLaN 数据集的领域内语言提示进行训练，但 OmniVLA 和 LeLaN 之间的性能差距在“行为”指标上尤为明显，该指标衡量的是遵循 OOD 语言提示的成功率。这一结果表明，利用一个更大规模的预训练 LLM 为导航任务提供了更强的先验知识。OmniVLA 的性能也优于 CounterfactualVLA，尽管 CounterfactualVLA 是用更多样化的语言提示进行训练的，但 OmniVLA 利用了更大的骨干网络和更大规模的训练数据组合，使其对 OOD 提示更具通用性和灵活性。我们还在相同环境下使用我们的提示评估了 NaVILA。然而，由于提示风格的领域差距，NaVILA 在所有指标上都失败了，得分为 0.0：它需要详细的、逐步的指令，例如“右转，直行，看到 X 时停止”，即使 X 已经可见且在附近，这与我们所追求的开放集合、自然语言提示格式不符。

图 4 展示了 OmniVLA 和基线模型在两个语言条件导航任务上的执行轨迹。这两个场景都特别具有挑战性，因为从机器人的初始位置看不到目标物体，这要求策略利用提示中的其他信息来到达目标。我们的方法成功地遵循了语言指令并到达了目标物体，而基线模型 LeLaN 和 CoW 则失败了，导航到了错误的物体。

\textbf{数据集消融实验。} 我们进行了一项消融研究，以突显在保持模型架构固定的情况下，使用更大、更多样化的数据混合进行训练的好处。我们训练了多个版本的 OmniVLA-edge，每个版本只使用单一模态，并将它们与我们的多模态 OmniVLA-edge 策略进行比较。表 III 总结了用于训练模型的数据集，并报告了每个任务的目标到达率。在此评估中，我们还引入了卫星图像作为一种目标模态，并在与 2D 目标位姿导航相同的环境和设置下评估其性能。如表 III 所示，OmniVLA 表现出明显的优势，尤其是在以语言和卫星图像为条件的导航中。通过在多种模态上联合训练，OmniVLA 从所有数据集中学习核心导航行为，并与在有限数据上训练的单模态策略相比，实现了对未见环境更强的泛化能力。此外，我们观察到，包含高度变化的跨实体数据，即 BDD-V 数据集，使得策略能够在大约一半没有 BDD-V 训练的策略失败的案例中取得成功，这突显了 OmniVLA 吸收和受益于多样化数据源的能力。

\subsubsection{使用 OmniVLA 进行全模态条件化}

通过在全模态任务表征上进行训练，OmniVLA 可以学习遵循多种目标信号。为回答 Q2，我们进行了实验，任务通过提供 2D 目标位姿（何处？）和行为性语言指令（如何？）来指定，实验在 10 个不同的环境中进行。在每个环境中，我们提供一个距离机器人初始位置 25-100 米的 2D 位置，同时附带一个语言提示，该提示指定了机器人为到达目标必须遵循的行为，例如“沿着墙移动”、“在草地上移动”以及“在物体 A 和 B 之间移动”。这些提示是分布外的（OOD），并未包含在训练数据集中。据我们所知，尚无现有工作探索过导航中如此具有挑战性的模态组合。

由于没有现有的方法可以同时处理这两种模态，我们将我们的方法与使用 2D 位姿条件训练的专用基线（NoMaD 和 MBRA-pose，表 IV）进行比较。OmniVLA 展示了其能够同时关注提示和目标位姿信息的能力，实现了 80\% 的成功率，并且语言能力相比表 II 仅下降了 5\%。较小的 OmniVLA 变体由于模态容量有限而无法处理语言指令。其他 VLA 基线（SmolVLA 和 MiniVLA）也无法解决这个任务。

\subsubsection{OmniVLA 对新目标模态的适应}

为评估 OmniVLA 作为导航“基础模型”的能力并回答 Q3，我们评估了三个方面：（i）学习一种新的目标模态，（ii）在新数据集上进行微调，以及（iii）控制新的机器人实体。

对于（i），我们建立了一个受控实验，在没有卫星模态的情况下预训练 OmniVLA-edge。然后，我们冻结模型的大部分参数，并将用于自目心图像目标的视觉编码器替换为用于卫星图像的编码器，仅训练被替换的编码器。通过这个模型，我们评估了我们的模型利用新目标模态的适应性。我们发现 OmniVLA-edge 确实能够适应这种新模态，相比仅在卫星目标导航上训练的专用策略，性能显著提升（0.19 $\rightarrow$ 0.62）。这表明 OmniVLA 学到的跨模态表征可以促进学习新任务，同时保留在预训练期间学到的有用导航能力。

为评估（ii），我们进行两个实验来研究 OmniVLA 适应新数据的能力：（a）使用小数据集适应新环境，以及（b）从一个新的可用数据集中学习新的语言领域。对于（a），我们从未在训练中见过的测试环境中收集了 1.2 小时的数据，包括 2D 目标位姿和卫星图像，并对训练好的 OmniVLA-edge 进行微调以评估其适应性。表 V 显示，微调后的策略在两种模态上的性能都有所提升，这表明 OmniVLA 可以有效地适应新环境。对于（b），我们使用最近发布的 CounterfactualVLA 数据集对我们的 OmniVLA 进行微调，该数据集包含更多样化的语言提示。为了稳定微调过程，我们采样批次以维持我们先前训练中使用的模态之间的平衡。具体来说，在微调 OmniVLA 时，一半先前从 LeLaN 采样的数据被 CounterfactualVLA 数据替换。微调后的策略在（SR, Behavior, SR\textsuperscript{S}, SR\textsuperscript{C}）上取得了 (0.825, 0.700, 1.000, 0.765) 的分数，优于原始分数 (0.725, 0.650, 1.000, 0.647)，展示了 OmniVLA 从新兴数据集中学习新语言领域的能力。由于 CounterfactualVLA 数据集为多样化的语言提示提供了原始和反事实的机器人动作，我们的策略学会了更好的避障和指令遵循能力，在“行为”和“SR\textsuperscript{C}”指标上取得了显著进步。

为展示我们的模型在多种机器人数据上训练后的通用性，对于（iii），我们在另外两个平台上进行了部署：轮式移动机器人 VizBot 和四足机器人 Go1。我们在机器人上安装了不同的摄像头，并在最具挑战性的语言条件导航任务上对它们进行了测试。如图 6 所示，我们最好的 OmniVLA 使两种机器人都能遵循语言指令并到达目标位置。

\subsection{结论}

在本文中，我们介绍了 OmniVLA，一个用于机器人导航的全模态视觉-语言-动作模型。我们的策略在一个统一的框架内，灵活地解读多种目标模态，包括语言提示、2D 位姿、自目心目标图像以及它们的组合。通过从预训练的 VLA 初始化，OmniVLA 利用了互联网规模的知识，并在超过 9,500 小时的、涵盖十个机器人平台的导航经验数据上进行训练。

广泛的真实世界评估表明，OmniVLA 持续优于先前的基线模型，在各种模态上都取得了更强的性能，展现了鲁棒的分布外泛化能力，并能够遵循多样的多模态任务。此外，我们展示了其基础模型的特性，包括对新模态（例如卫星图像）的适应性、使用额外数据进行高效微调，以及跨不同机器人实体的迁移能力。

我们相信，OmniVLA 代表了朝着实现广泛通用的导航策略迈出的重要一步。为进一步推动这一方向的发展，我们将发布我们的模型和训练代码，作为开发可扩展且灵活的机器人导航模型的资源。特别是在以语言为条件的导航方面，无论是在目标到达还是指令遵循上，仍有提升的空间。该领域的进展将极大地受益于更大规模、更精心策划的数据集，我们希望社区能致力于此。
%%%%%%%%%%%%%%%%%%%%%%%%%%%%%%%%%%%%%%



%%%%%%%%%%%%%%%%%%%%%%%%%%%%%%%%%%%%%%%%%%
%%%%%%%%%%%%%%%%%%%%%%%%%%%%%%%%%%%%%%%%%%

\refstepcounter{section}
\section*{译文题目：GSPO 组序列策略优化}
\addcontentsline{toc}{section}{\protect\numberline{\thesection}译文题目：GSPO 组序列策略优化}
{\zihao{3}
\noindent \textbf{原文题目：Group Sequence Policy Optimization} \\
\textbf{原文来源：} Zheng C, Liu S, Li M, et al. Group Sequence Policy Optimization[J]. arXiv preprint arXiv:2507.18071, 2025. \\
}

\vskip 5mm

 {
\setlength{\baselineskip}{16pt}\selectfont{
\noindent {\heiti 摘\quad 要\quad }
本文介绍了组序列策略优化（Group Sequence Policy Optimization, GSPO），一种我们为训练大型语言模型而设计的稳定、高效且高性能的强化学习算法。与先前采用词元级别（token-level）重要性比率的算法不同，GSPO 基于序列似然（sequence likelihood）来定义重要性比率，并执行序列级别（sequence-level）的裁剪、奖励和优化。我们证明，与 GRPO 算法相比，GSPO 实现了卓越的训练效率和性能，显著稳定了混合专家（Mixture-of-Experts, MoE）模型的强化学习训练，并有潜力简化强化学习基础设施的设计。GSPO 的这些优点为最新的 Qwen3 模型带来了显著的性能提升。
\par}}


% %  {\noindent
% % {\heiti 关键词 \quad }{*关键词（中文关键字与英文关键字对应且一致，应有5-7个关键词)；关键词；关键词*  }
% }
% \zihao{5-}{\heiti 中图法分类号\quad } TP\rm{\quad \quad \quad     }
% {\heiti DOI号:\quad } *投稿时不提供DOI号


% \vskip 7mm


\subsection{引言}
强化学习（Reinforcement learning, RL）已成为一个扩展语言模型的关键范式 (OpenAI, 2024; DeepSeek-AI, 2025; Qwen, 2025b;a)。通过大规模强化学习，语言模型能够通过更深、更长的推理过程，培养解决复杂问题的能力，例如竞赛级别的数学和编程问题。

为了通过更大的计算投入成功地扩展强化学习，首要前提是维持稳定和鲁棒的训练动态。然而，当前最先进的强化学习算法，以 GRPO (Shao et al., 2024) 为例，在训练巨型语言模型时表现出严重的稳定性问题，常导致灾难性且不可逆的模型崩溃 (Qwen, 2025a; MiniMax, 2025)。这种不稳定性阻碍了通过持续的强化学习来推动语言模型能力边界的努力。

在本文中，我们发现 GRPO 的不稳定性源于其算法设计中对重要性采样权重的根本性误用与失效。这引入了高方差的训练噪声，该噪声随着响应长度的增加而逐步累积，并被裁剪机制进一步放大，最终导致模型崩溃。

为解决这些核心限制，我们提出了组序列策略优化（Group Sequence Policy Optimization, GSPO），一种用于训练大型语言模型的新型强化学习算法。GSPO 的关键创新在于其基于序列似然的重要性比率定义具有坚实的理论基础 (Zheng et al., 2023)，这与重要性采样的基本原则相符。此外，GSPO 将归一化奖励计算为对同一查询的多个响应的优势（advantages），确保了序列级别的奖励与优化之间的一致性。

\subsection{预备知识}
\textbf{符号表示} \quad 在本文中，一个由 $\theta$ 参数化的自回归语言模型被定义为一个策略 $\pi_\theta$。我们用 $x$ 表示查询，用 $\mathcal{D}$ 表示查询集。给定对查询 $x$ 的一个响应 $y$，其在策略 $\pi_\theta$ 下的似然表示为 $\pi_\theta(y|x) = \prod_{t=1}^{|y|} \pi_\theta(y_t|x, y_{<t})$，其中 $|y|$ 表示 $y$ 中的词元数量。

\textbf{近端策略优化 (PPO)} \quad PPO (Schulman et al., 2017) 使用从旧策略 $\pi_{\theta_{\text{old}}}$ 生成的样本，通过裁剪机制将策略更新限制在旧策略的一个邻近区域内。
\textbf{组相对策略优化 (GRPO)} \quad GRPO (Shao et al., 2024) 通过计算同一查询的一组响应中每个响应的相对优势，从而绕过了对价值模型的需求。

\subsection{动机}
模型规模、稀疏性（例如，在混合专家模型中）和响应长度的增长，要求在强化学习期间使用大的 rollout 批次大小（rollout batch size）以最大化硬件利用率。为了提高样本效率，标准做法是将一个大的 rollout 数据批次划分为多个小批次（mini-batches）进行梯度更新。

尽管像裁剪这样的机制旨在管理这种离策略差异，但我们发现 GRPO 中存在一个更根本性的问题：其目标函数是病态的（ill-posed）。当在长响应任务上训练大型模型时，这个问题变得尤为突出，导致灾难性的模型崩溃。GRPO 目标函数的病态性质源于对重要性采样权重的误用。

\subsection{算法}
\textbf{GSPO：组序列策略优化} \quad 基于这一直接的观察，我们提出了组序列策略优化（Group Sequence Policy Optimization, GSPO）算法。GSPO 采用以下序列级别的优化目标：
\begin{equation}
J_{\text{GSPO}}(\theta) = \mathbb{E}_{x \sim \mathcal{D}, \{y_i\}_{i=1}^G \sim \pi_{\theta_{\text{old}}}(\cdot|x)} \left[ \frac{1}{G} \sum_i \min \left( s_i(\theta) \hat{A}_i, \text{clip}(s_i(\theta), 1-\epsilon, 1+\epsilon) \hat{A}_i \right) \right]
\end{equation}
其中我们采用基于组的优势估计，并基于序列似然定义重要性比率 $s_i(\theta)$。

\textbf{梯度分析} \quad 我们可以如下推导 GSPO 目标的梯度。相比之下，GSPO 和 GRPO 之间的根本区别在于它们如何对词元对数似然的梯度进行加权。在 GRPO 中，词元根据其各自的“重要性权重”进行加权。

\subsubsection{GSPO-token：一种词元级别的目标函数变体}

在像多轮强化学习这样的场景中，我们可能希望进行比序列级别更细粒度的优势调整。为此，我们引入了 GSPO 的一个词元级别的目标函数变体，即 GSPO-token，以允许进行逐词元的优势定制：
\begin{equation}
J_{\text{GSPO-token}}(\theta) = \mathbb{E}_{x \sim \mathcal{D}, \{y_i\}_{i=1}^G \sim \pi_{\theta_{\text{old}}}(\cdot|x)} \left[ \frac{1}{G} \sum_i \frac{1}{|y_i|} \sum_t \min \left( s_{i,t}(\theta) \hat{A}_{i,t}, \text{clip}(s_{i,t}(\theta), 1-\epsilon, 1+\epsilon) \hat{A}_{i,t} \right) \right]
\end{equation}
其中
\begin{equation}
s_{i,t}(\theta) = \text{sg}[s_i(\theta)] \cdot \frac{\pi_\theta(y_{i,t}|x, y_{i,<t})}{\text{sg}[\pi_\theta(y_{i,t}|x, y_{i,<t})]},
\end{equation}
$\text{sg}[\cdot]$ 表示只取数值但停止梯度，对应于 PyTorch 中的 detach 操作。GSPO-token 的梯度可以推导为：
\begin{align}
\nabla_\theta J_{\text{GSPO-token}}(\theta) &= \nabla_\theta \mathbb{E}_{x \sim \mathcal{D}, \{y_i\}_{i=1}^G \sim \pi_{\theta_{\text{old}}}(\cdot|x)} \left[ \frac{1}{G} \sum_i \frac{1}{|y_i|} \sum_t s_{i,t}(\theta) \hat{A}_{i,t} \right] \\
 &= \mathbb{E}_{x \sim \mathcal{D}, \{y_i\}_{i=1}^G \sim \pi_{\theta_{\text{old}}}(\cdot|x)} \left[ \frac{1}{G} \sum_i \text{sg}[s_i(\theta)] \frac{1}{|y_i|} \sum_t \hat{A}_{i,t} \nabla_\theta \frac{\pi_\theta(y_{i,t}|x, y_{i,<t})}{\text{sg}[\pi_\theta(y_{i,t}|x, y_{i,<t})]} \right] \\
 &= \mathbb{E}_{x \sim \mathcal{D}, \{y_i\}_{i=1}^G \sim \pi_{\theta_{\text{old}}}(\cdot|x)} \left[ \frac{1}{G} \sum_i \frac{\pi_\theta(y_i|x)}{\pi_{\theta_{\text{old}}}(y_i|x)} \frac{1}{|y_i|} \sum_t \hat{A}_{i,t} \nabla_\theta \log \pi_\theta(y_{i,t}|x, y_{i,<t}) \right]
\end{align}
注意，项 $\frac{\pi_\theta(y_{i,t}|x, y_{i,<t})}{\text{sg}[\pi_\theta(y_{i,t}|x, y_{i,<t})]}$ 的数值为 1，因此 $s_{i,t}(\theta)$ 在数值上等于 $s_i(\theta)$。比较公式 (5) 和 (13)，以及公式 (10) 和 (17)，当我们将响应 $y_i$ 中所有词元的优势设置为相同的值时（即 $\hat{A}_{i,t} = \hat{A}_i$），GSPO-token 和 GSPO 在优化目标、裁剪条件和理论梯度上是数值等价的，而 GSPO-token 享有调整每个词元优势的更高灵活性。

\subsubsection{实验与讨论}

\textbf{实证结果}

我们使用一个从 Qwen3-30B-A3B-Base 冷启动微调的模型进行实验，并报告了在 AIME'24（32个样本的平均 Pass@1）、LiveCodeBench（202410-202502，8个样本的平均 Pass@1）和 CodeForces（Elo 等级分）基准测试上的训练奖励曲线以及模型性能曲线。在强化学习训练期间，每个批次的 rollout 数据被划分为四个小批次（mini-batches）进行梯度更新。在 GSPO 中，我们将公式(5)中的左右裁剪范围分别设置为 3e-4 和 4e-4。我们以 GRPO 作为基线进行比较，并将其在公式(2)中的左右裁剪范围分别设置为 0.2 和 0.27，我们对这些值进行了仔细调整以确保公平比较。

请注意，GRPO 需要路由回放（Routing Replay）训练策略来实现 MoE 强化学习的正常收敛，我们将在 \S 5.3 节中进一步讨论，而 GSPO 则无需此策略。

图  显示，使用 GSPO 的训练过程全程保持稳定。我们观察到，通过增加训练计算量、定期更新查询集和延长生成长度，GSPO 能够实现持续的性能提升。此外，GSPO 在训练效率上也优于 GRPO，在相同的训练计算量和消耗的查询下，实现了更好的训练准确率和基准性能。最后，我们已成功将 GSPO 应用于最新的 Qwen3 模型的强化学习训练中，有力地证明了 GSPO 在释放大规模语言模型强化学习扩展能力方面的有效性。

\textbf{关于裁剪比例的奇特观察}

GSPO 与 GRPO 的一个关键区别在于其裁剪整个响应而非单个词元的做法。特别地，如图所示，我们观察到 GSPO 和 GRPO 之间被裁剪词元比例存在两个数量级的差异（调整裁剪范围并不会改变这种量级上的差异）。然而，尽管裁剪了显著更多的词元，从而用于训练（或梯度估计）的词元更少，GSPO 仍然比 GRPO 实现了更高的训练效率。这一反直觉的发现——即裁剪更大部分的词元反而带来了更优的训练效率——进一步表明 GRPO 的词元级别梯度估计本质上是充满噪声且在样本利用上效率低下。相比之下，GSPO 的序列级别方法提供了一个更可靠和有效的学习信号。

\textbf{GSPO 对 MoE 训练的裨益}

\textbf{背景} \quad 与密集模型的强化学习训练相比，MoE 模型的稀疏激活特性引入了独特的稳定性挑战。我们特别发现，在采用 GRPO 算法时，MoE 模型的专家激活波动性（expert-activation volatility）会阻碍强化学习训练的正常收敛。具体来说，经过一次或多次梯度更新后，针对同一响应激活的专家可能会发生显著变化。例如，对于48层的 Qwen3-30B-A3B-Base 模型，在每次强化学习梯度更新后，对于相同的 rollout 样本，在新策略 $\pi_\theta$ 下激活的专家中，约有10\%与旧策略 $\pi_{\theta_{\text{old}}}$ 下激活的专家不同。这种现象在更深的 MoE 模型中更为突出，导致词元级别的重要性比率 $w_{i,t}(\theta) = \frac{\pi_\theta(y_{i,t}|x, y_{i,<t})}{\pi_{\theta_{\text{old}}}(y_{i,t}|x, y_{i,<t})}$ 发生剧烈波动，并进一步失效，正如在第3节和第4.2节中讨论的，从而阻碍了强化学习训练的正常收敛。

\textbf{我们先前的方法} \quad 为应对这一挑战，我们先前采用了路由回放（Routing Replay）训练策略。具体来说，我们缓存 $\pi_{\theta_{\text{old}}}$ 中被激活的专家，并在计算重要性比率 $w_{i,t}(\theta) = \frac{\pi_\theta(y_{i,t}|x, y_{i,<t})}{\pi_{\theta_{\text{old}}}(y_{i,t}|x, y_{i,<t})}$ 时，在 $\pi_\theta$ 中“回放”这些路由模式。通过这种方式，对于每个词元 $y_{i,t}$，$\pi_\theta(y_{i,t}|x, y_{i,<t})$ 和 $\pi_{\theta_{\text{old}}}(y_{i,t}|x, y_{i,<t})$ 共享相同的激活网络，因此我们可以恢复词元级别重要性比率的稳定性，并确保在梯度更新中对一致的激活网络进行优化。图  表明，路由回放是 GRPO 训练 MoE 模型正常收敛的一项关键技术。

\textbf{GSPO 的优势} \quad 尽管路由回放能使 MoE 模型的 GRPO 训练正常收敛，但其复用路由模式的做法会带来额外的内存和通信开销，并且也可能限制 MoE 模型的实际能力。相比之下，如图  所示，GSPO 消除了对路由回放的依赖，完全能够常规地计算重要性比率 $s_i(\theta)$，并实现正常收敛和稳定优化。其关键洞见在于，GSPO 仅关注序列似然（即 $\pi_\theta(y_i|x)$），而对单个词元的似然（即 $\pi_\theta(y_{i,t}|x, y_{i,<t})$）不敏感。由于 MoE 模型始终保持其语言建模能力，序列似然不会剧烈波动。总之，GSPO 从根本上解决了 MoE 模型中的专家激活波动性问题，从而无需像路由回放这样的复杂变通方法。这不仅简化和稳定了训练过程，还允许模型在没有人为限制的情况下充分发挥其全部能力。

\textbf{GSPO 对强化学习基础设施的裨益}

鉴于训练引擎（如 Megatron）和推理引擎（如 SGLang 和 vLLM）之间的精度差异，在实践中，我们通常使用训练引擎在旧策略 $\pi_{\theta_{\text{old}}}$ 下重新计算采样响应的似然。然而，GSPO 仅使用序列级别而非词元级别的似然进行优化，直观上，前者对精度差异的容忍度要高得多。因此，GSPO 使得直接使用推理引擎返回的似然进行优化成为可能，从而避免了使用训练引擎进行重新计算的需要。这在部分 rollout（partial rollout）和多轮强化学习等场景以及训练与推理分离的框架中尤其有益。



\subsubsection{结论}
我们提出了组序列策略优化（Group Sequence Policy Optimization, GSPO），一种用于训练大型语言模型的新的强化学习算法。遵循重要性采样的基本原则，GSPO 基于序列似然定义重要性比率，并执行序列级别的裁剪、奖励和优化。与 GRPO 相比，GSPO 在训练稳定性、效率和性能方面表现出显著的优越性，并对大规模 MoE 模型的强化学习训练展现出特别的功效。

\begingroup
    \linespreadsingle{}
    \printbibliography[title={外文翻译参考文献}]
\endgroup
